\section{P\'olya urns and beta prediction}
\label{sec:annex-p15-polya-urns}

\begin{annexitem}{annex:p15.polya-construction}{thm:p15-polya-construction}{Recursive draws and the bounded martingale}
\begin{proof}
The independent uniform sequence exists by the countable product
construction. At step \(n+1\), all earlier draws and \(X_n\) are
measurable functions of \(U_1,\ldots,U_n\).
The new uniform is independent of this sigma-field. Integrating its
indicator at each earlier outcome gives conditional probability \(X_n\).
The weights remain positive, their sum is \(g+r+cn\), and the recursion
therefore defines every draw and every fraction.
This also proves the conditional assertion on the smaller draw
sigma-field \(\mathcal F_n\).

Since \(S_{n+1}=S_n+B_{n+1}\),
\[
 \frac{\alpha+S_n+B_{n+1}}{t+n+1}
       -\frac{\alpha+S_n}{t+n}
             =\frac{B_{n+1}-X_n}{t+n+1}.
\]
Conditional expectation makes this increment have mean zero.
Boundedness gives integrability and \(EX_n=X_0=\alpha/t\).
The nonnegative martingale convergence theorem gives a finite
almost-sure limit in \([0,1]\).
For every fixed \(p>0\), \(|X_n-X_\infty|^p\leq1\);
dominated convergence gives convergence of these moments to zero,
which is the stated \(L^p\) convergence.

For \(n\geq1\), direct subtraction gives
\[
 X_n-\frac{S_n}{n}
       =\frac{\alpha-t(S_n/n)}{t+n},\qquad
 \left|X_n-\frac{S_n}{n}\right|
             \leq\frac{\max(\alpha,\beta)}{t+n},
\]
because \(S_n/n\in[0,1]\). The deterministic bound tends to zero.
\end{proof}
\end{annexitem}

\begin{annexitem}{annex:p15.polya-word-law}{thm:p15-polya-word-law}{The source's permuted numerator argument}
\begin{proof}
Along a word with \(s\) green draws, the green choices have successive
numerators \(\alpha,\alpha+1,\ldots,\alpha+s-1\);
the red choices have successive numerators
\(\beta,\ldots,\beta+n-s-1\).
All draws have successive denominators \(t,t+1,\ldots,t+n-1\).
Multiplying their conditional probabilities gives the stated ratio,
regardless of how the two numerator lists are interleaved.
Every finite prefix has positive probability because both parameters
are positive, so this multiplication uses well-defined conditional
probabilities at every prefix. The length-zero case has probability one.

There are \(\binom ns\) distinct binary words with count \(s\).
Their events are disjoint and exhaust \(\{S_n=s\}\); summing gives
the count law. The probabilities sum to one because those events
partition the sample space, not because a separate binomial identity
has been assumed.

For a permutation moving finitely many coordinates, choose \(N\)
containing its support and the coordinates of a specified cylinder.
The joint probability of every length-\(N\) word depends only on its
count, which the permutation preserves. Summing over unspecified
coordinates gives equality on the specified cylinder.
The cylinder pi-system generates the full sequence sigma-field,
so uniqueness of probability measures gives exchangeability on all
measurable sequence events.
\end{proof}
\end{annexitem}

\begin{annexitem}{annex:p15.polya-mixture-representation}{thm:p15-polya-mixture-representation}{An explicit mixture and its almost-sure frequency}
\begin{proof}
The beta normalization and gamma identity are already established in
\cref{thm:p15-beta-gamma-convolution}.
The beta variable and the independent uniform sequence can therefore
be realized on the product of their probability spaces.
Each map \((\theta,v)\mapsto\mathbf1_{\{v_j\leq\theta\}}\)
is measurable, so all sequence events used below are measurable.

For fixed \(0<\theta<1\), independence of the uniforms gives
probability \(\theta^s(1-\theta)^{n-s}\) to a specified word.
Multiplying this by the beta density and integrating over \(D\)
gives the joint integral in the statement, by nonnegative product
integration. The beta law gives its endpoints mass zero, and all
displayed endpoint density values are assigned zero.

Integrating over all parameters yields
\[
 \frac{B(\alpha+s,\beta+n-s)}{B(\alpha,\beta)}
           =\frac{(\alpha)_s(\beta)_{n-s}}{(t)_n}.
\]
The equality follows by the beta-gamma identity and the repeated
gamma recursion
\(\Gamma(a+k)/\Gamma(a)=(a)_k\), including \(k=0\).
This is the urn word probability. Equality on every finite prefix,
then on every finite-coordinate cylinder by marginalization, gives
equality of the two probability measures on
\(\{0,1\}^{\mathbb N}\) by the pi-lambda theorem.

For fixed \(\theta\in[0,1]\), the section sequence
\(\mathbf1_{\{V_j\leq\theta\}}\) is i.i.d. Bernoulli of mean \(\theta\).
The strong law \cref{thm:p15-iid-series-strong-law} gives convergence
of its averages to \(\theta\) with uniform-sequence probability one.
The event of that convergence is the countable measurable event
\[
 \bigcap_{k\geq1}\ \bigcup_{N\geq1}\ \bigcap_{n\geq N}
 \left\{\left|n^{-1}\sum_{j=1}^nC_j-\Theta\right|<1/k\right\}.
\]
Every parameter section has measure one. Tonelli applied to its
indicator gives product probability one. Thus no single null subset
of the uniform space is presumed to work for every parameter.
\end{proof}
\end{annexitem}

\begin{annexitem}{annex:p15.polya-beta-limit}{thm:p15-polya-beta-limit}{Transferring the joint limit through the full sequence law}
\begin{proof}
For a binary sequence \(b\), define \(L(b)\) to be its limiting
frequency if that limit exists and zero otherwise.
Existence is a measurable Cauchy condition on the coordinate averages.
The limsup is measurable and agrees with the limit on that event,
so the piecewise map \(L\) is measurable.

The construction theorem gives \(X_\infty=L(B)\) almost surely.
The mixture theorem gives \(\Theta=L(C)\) almost surely.
The two complete sequence laws are equal. Their pushforwards under
the measurable map \(b\mapsto(b,L(b))\) are therefore equal.
Consequently \((B,X_\infty)\) and \((C,\Theta)\) have the same
joint law. This proves both the beta limit and its joint finite-word
identity.

For completeness, the conditional Bernoulli assertion means that
for each finite word \(b\) with count \(s\),
\[
 P(B_1=b_1,\ldots,B_n=b_n\mid\sigma(X_\infty))
                =X_\infty^s(1-X_\infty)^{n-s}
                         \quad\text{almost surely}.
\]
At endpoints, an exponent zero is assigned the value one in this
finite Bernoulli product; the limit law gives the endpoints mass zero
in the positive-parameter case. Integrating this proposed conditional
probability on \(\{X_\infty\in D\}\) gives exactly the established
joint integral. Such events generate \(\sigma(X_\infty)\), so the
all-event conditional identity follows. There are countably many
finite words, permitting one common almost-sure event for all their
identities. Marginalization gives the assertion for arbitrary finite
coordinate subfamilies.
\end{proof}
\end{annexitem}

\begin{annexitem}{annex:p15.polya-posterior-prediction}{thm:p15-polya-posterior-prediction}{The normalized update and full predictive count}
\begin{proof}
Every observed word has positive probability.
Divide its joint parameter integral by its word probability.
The resulting conditional density is
\[
       \frac{\theta^{a-1}(1-\theta)^{b-1}}{B(a,b)}
                         \quad(0<\theta<1),
\]
with \(a=\alpha+s,b=\beta+n-s\).
The beta mean and variance from
\cref{thm:p15-beta-gamma-convolution} give \(a/T=X_n\) and
\(ab/[T^2(T+1)]=X_n(1-X_n)/(t+n+1)\).
Since \(\mathcal F_n\) consists of unions of these finitely many word
events, the formulas define the stated conditional variables on all
of \(\mathcal F_n\).

For a specified future word of length \(m\) with \(k\) green draws,
divide the joint length-\(n+m\) probability by the past word probability.
Cancellation of rising products gives
\((a)_k(b)_{m-k}/(T)_m\).
There are \(\binom mk\) such future words, which gives the count law.
Equivalently it follows by integrating the Bernoulli word probability
against the updated beta density.

A single future indicator has conditional mean \(a/T\) and variance
\(ab/T^2\). Two distinct future indicators have conditional joint
success probability
\[
                        \frac{a(a+1)}{T(T+1)}.
\]
Subtracting \((a/T)^2\) gives conditional covariance
\(ab/[T^2(T+1)]\). Adding all \(m\) variances and the \(m(m-1)\)
ordered covariance terms yields
\[
 \frac{mab}{T^2}+\frac{m(m-1)ab}{T^2(T+1)}
       =\frac{mab(T+m)}{T^2(T+1)}.
\]
The variance exceeds the fixed-probability binomial variance when
\(m>1\), since \(a,b>0\) and \((T+m)/(T+1)>1\).
\end{proof}
\end{annexitem}

\begin{annexitem}{annex:p15.polya-moments}{thm:p15-polya-moments}{Mixed beta moments and the common-parameter covariance}
\begin{proof}
Integrating the mixed power against the beta density gives
\(B(\alpha+k,\beta+l)/B(\alpha,\beta)\).
The same gamma recursion as in the mixture proof gives the displayed
rising-product ratio. Setting \((k,l)=(1,0),(2,0)\) gives
\[
 EX_\infty=\alpha/t,\qquad
 EX_\infty^2=\alpha(\alpha+1)/(t(t+1)),
\]
and subtraction gives the stated variance.

Exchangeability gives \(EB_i=\alpha/t\) at every coordinate.
The common-parameter finite joint identity gives
\(E(B_iB_j)=EX_\infty^2\) for distinct coordinates.
Thus their covariance equals \(\operatorname{Var}(X_\infty)\).
A Bernoulli variable of probability \(\alpha/t\) has variance
\(\alpha\beta/t^2>0\); dividing the covariance by this variance
gives the correlation \(1/(t+1)\).

Finite linearity gives \(ES_n=n\alpha/t\).
Adding variances and ordered cross terms gives
\[
 \operatorname{Var}(S_n)
   =\frac{n\alpha\beta}{t^2}
       +\frac{n(n-1)\alpha\beta}{t^2(t+1)}
   =\frac{n\alpha\beta(t+n)}{t^2(t+1)}.
\]
Because \(X_n=(\alpha+S_n)/(t+n)\), division by \((t+n)^2\)
gives its variance. Both means are \(\alpha/t\).

The posterior identity gives
\(E(X_\infty X_n)=E(X_n E(X_\infty\mid\mathcal F_n))=EX_n^2\).
All variables are bounded, so the products are integrable and
conditional pullout is valid. Therefore
\[
 E(X_\infty-X_n)^2
        =EX_\infty^2-EX_n^2
        =\operatorname{Var}(X_\infty)-\operatorname{Var}(X_n)
        =\frac{\alpha\beta}{t(t+1)(t+n)}.
\]
At \(n=0\) the formulas give \(S_0=0\), constant \(X_0=\alpha/t\)
and the correct initial error \(\operatorname{Var}(X_\infty)\).
\end{proof}
\end{annexitem}

\begin{annexitem}{annex:p15.polya-source-special-cases}{thm:p15-polya-source-special-cases}{The two source finite-mass calculations}
\begin{proof}
For \(\alpha=\beta=1\), rising products are factorials, so
\[
 P(S_n=s)=\frac{n!}{s!(n-s)!}\frac{s!(n-s)!}{(n+1)!}
                           =\frac1{n+1}.
\]
The fraction is \((s+1)/(n+2)\).
For a continuous function \(f\) on \([0,1]\),
\((n+1)^{-1}\sum_{s=0}^n f((s+1)/(n+2))\) tends to
\(\int_0^1f(u)\,du\), by uniform continuity and the equal-mesh
Riemann sums. Almost-sure bounded convergence of \(X_n\) gives the
same limit as \(Ef(X_\infty)\). Bounded continuous tests determine
the law, proving the source's uniform limit directly.

For \(\alpha=2,\beta=1\),
\((2)_s=(s+1)!\), \((1)_{n-s}=(n-s)!\), and
\((3)_n=(n+2)!/2\). Multiplication by \(\binom ns\) gives
\(2(s+1)/((n+1)(n+2))\).
Summing the arithmetic sequence from \(s=0\) through \(k\) gives
the stated CDF for \(S_n\).

Now \(X_n=(S_n+2)/(n+3)\).
For fixed \(0<x<1\), its CDF is obtained by taking
\(k=\lfloor(n+3)x-2\rfloor\) and truncating to the allowed range.
For all sufficiently large \(n\), \(0\leq k<n\), and \(k/n\to x\).
The displayed count CDF tends to \(x^2\); at zero it is zero and
at one it is one. Thus the finite laws converge to the law with
CDF \(x^2\) on \([0,1]\), which has density \(2x\) on \((0,1)\).
The almost-sure fraction limit has this law, either by the
bounded-continuous characterization of weak convergence or the
already proved general beta limit.

Finally, for \(s_n/n\to x\),
\[
 nP(G_n=s_n+2)
       =\frac{2(s_n/n+1/n)}
               {(1+1/n)(1+2/n)}\longrightarrow2x.
\]
The unscaled mass is bounded by \(2/(n+2)\), hence tends to zero.
This explicitly retains the mesh factor between mass and density.
\end{proof}
\end{annexitem}

\begin{annexitem}{annex:p15.polya-parameter-boundaries}{thm:p15-polya-parameter-boundaries}{Small and large reinforcement and zero-weight processes}
\begin{proof}
Common scaling leaves \(\alpha,\beta\) unchanged, hence leaves every
word probability unchanged. The full sequence law is determined by
those probabilities.

For fixed \(g,r>0\), the beta parameter sum is \(t=(g+r)/c\),
and \(\alpha=pt,\beta=(1-p)t\).
The mean is \(p\), and the variance is \(p(1-p)/(t+1)\).
As \(c\downarrow0\), \(t\to\infty\). For every \(\epsilon>0\),
Chebyshev gives mass outside \((p-\epsilon,p+\epsilon)\) at most
\(p(1-p)/((t+1)\epsilon^2)\), which tends to zero.
For a bounded continuous test on \([0,1]\), uniform continuity
and that bound show its integral tends to its value at \(p\);
this is weak convergence to \(\delta_p\).

As \(c\to\infty\), \(t\downarrow0\), and the mixed moment is
\[
          E[\Theta(1-\Theta)]=\frac{p(1-p)t}{t+1}\longrightarrow0,
          \qquad E\Theta=p.
\]
For \(0<\delta<1/2\), the mass of \([\delta,1-\delta]\) is at most
this mixed moment divided by \(\delta(1-\delta)\).
Let \(f\) be continuous on \([0,1]\), and subtract its endpoint
interpolant \(\ell(x)=(1-x)f(0)+xf(1)\).
The difference \(f-\ell\) is continuous and zero at both endpoints.
Choose \(\delta\) so its absolute value on the two endpoint
neighborhoods is arbitrarily small; its integral on the middle
interval tends to zero by the bound.
But \(E\ell(\Theta)=(1-p)f(0)+pf(1)\) for every \(c\).
This proves the asserted endpoint-mixture weak limit.

At exactly \(c=0\), all weights stay constant.
The independent-uniform construction now uses the same threshold
\(p\) at every step, so the indicators are independent Bernoulli
variables of that probability. Their strong law gives the frequency
limit. If an initial color has zero weight, its proportional chance
is zero and remains zero after every draw; induction gives the
degenerate process and limit. Positive total initial weight keeps
every denominator positive. With zero total initial weight there
is no proportional probability for the first draw.
\end{proof}
\end{annexitem}

\begin{annexitem}{annex:p15.polya-worked-prediction}{thm:p15-polya-worked-prediction}{Prediction and model comparison with distinct conditionings}
\begin{proof}
Here \(\alpha=2,\beta=1,t=3\). The ordered word has probability
\((2)_2(1)_1/(3)_3=6/60=1/10\).
The three words with count two each have this probability, so the
count event has probability \(3/10\).
The updated parameters are \(a=4,b=2,T=6\).
The beta mean is \(4/6=2/3\), and its variance is
\(8/(36\cdot7)=2/63\).

For the next two draws, the denominator \((6)_2\) is 42.
The three count numerators are \((2)_2=6\),
\(2(4)_1(2)_1=16\), and \((4)_2=20\).
They give \(1/7,8/21,10/21\), which sum to one.
The variance formula gives
\(2\cdot4\cdot2\cdot8/(36\cdot7)=32/63\).
The independent plug-in joint success probability is
\((2/3)^2=4/9\), different from \(10/21\).

Under independent fair draws the specified ordered data have probability
\(1/8\), so the ratio of data probabilities is
\((1/10)/(1/8)=4/5\).
To give the odds assertion its explicit probability space, first choose
a model indicator with probabilities \(\pi_U,\pi_F>0\), summing to one,
then draw the sequence from its indicated law. Bayes' finite partition
formula gives posterior odds
\[
       \frac{\pi_U(1/10)}{\pi_F(1/8)}
                    =\frac45\,\frac{\pi_U}{\pi_F}.
\]
With \(\pi_U=\pi_F=1/2\), normalization gives \(4/(4+5)=4/9\).
This is a probability for the model indicator. It is not the urn
model's conditional success probability \(2/3\), which conditions
on both the model and the observed word.
\end{proof}
\end{annexitem}

